\begin{abstract}
Drug--target affinity (DTA) models often excel within a benchmark yet fail under distribution shift---across datasets, protein families, or structurally unconventional targets. This brittleness reflects a limitation of pair-based architectures: they memorize dataset-specific associations rather than learning transferable determinants of binding. We introduce Anchor Transfer Learning, which reformulates DTA prediction as a comparison problem. Instead of scoring a protein--drug pair in isolation, the model conditions each prediction on an anchor protein already known to bind a chemically similar drug, shifting the question from ``does protein $P$ bind drug $D$?'' to ``how does $P$ compare with a known binder of a related compound?'' At test time, anchors are retrieved from the training set by Tanimoto chemical similarity after excluding canonicalized chemical duplicates, requiring no privileged information about the evaluation data. Under a cross-dataset protocol with verified zero canonical compound overlap, V2-650M achieves per-protein CI 0.642 and AUROC 0.669 on the Davis kinase benchmark (442 proteins, 68 drugs), while DeepDTA falls to near-random performance (CI 0.521). Performance stratified by retrieval similarity shows that the model generalizes even to chemically dissimilar compounds (Tanimoto~$< 0.6$: CI 0.656, AUROC 0.651). Stratifying by anchor strength reveals that the 650M encoder compensates for weak anchors, while strong anchors equalize smaller encoders. A cosine similarity ablation confirms that the model learns drug-specific binding patterns beyond raw protein similarity, with 74\% of predictive variance orthogonal to embedding cosine distance. These results establish anchor-based transfer as a practical route to DTA models that generalize across datasets.
\end{abstract}
